\scanpage{175}
\begin{center}- 27 -\end{center}\hfill IX

b) Soit $n>0$ premier à l'exposant caractéristique de $k(s)$
et aux rangs des quotients des fibres géométriques de $K$ en $s$ et $t$
par leurs composantes neutres réduites, alors les deux membres de
(3.1.1) s'obtiennent en élevant $n$ respectivement à la puissance
d'exposant l'un et l'autre membre de (3.1.2), comme il résulte de la
structure connue des groupes algébriques commutatifs. Donc b) résulte de a).

c) Comme l'hypothèse $(K_{\bar t})^{\unsure{o|0}}_s$ unipotent équivaut aux relations
$\mu(t)=\alpha(t)=0$, et de même pour $K_s$, l'assertion c) résulte formellement
de b).
\figure{\begin{tikzpicture}\draw[line width=.6pt] (0,.13)--(.18,.14);\draw[line width=.6pt] (0,0)--(.18,-.02);\end{tikzpicture}}

d) L'hypothèse que $u_t$ soit une isogénie implique que $K_t$ est
fini, donc on est sous les conditions de c), et par suite $(K_{\bar s})^{\unsure{o|0}}_{\mathrm{red}}$ est
unipotent. Comme $G_s$ est de rang unipotent nul, on en déduit que $K_s$
est fini. D'ailleurs le fait que $u_t$ est une isogénie implique
$\dim G_t=\dim H_t$, et le fait que $G$ et $H$ sont plats localement de présentation finie implique $\dim G_s=\dim G_t$, $\dim H_s=\dim H_t$ (SGA 3 $\mathrm{VI_B}$ 4.3),
d'où $\dim G_s=\dim H_s$. Comme $\dim\operatorname{Ker}u_s=0$, on en conclut que
$\dim H_s=\dim u_s(G_s)$ (SGA 3 $\mathrm{VI_B}$ 1.2.), ce qui équivaut à $u_s$ plat
grâce à l'hypothèse $H_s$ lisse (SGA 3 $\mathrm{VI_A}$ 5.6). On en conclut qu'il
existe un voisinage ouvert $U$ de $s$ tel que $u$ soit plat et quasi-fini
en les points de la section unité sur $U$ (SGA 3 $\mathrm{VI_B}$ 2.2 (i) et 2.5 (ii)).
Il en résulte, comme pour $u_s$, que $u_{s'}$ est plat pour tout $s'\in U$, donc
$u|G_U$ est plat et quasi-fini, l'étant sur chaque fibre.

e) Pour la premier assertion, on se ramène aussitôt par
changement de base au cas où $S$ est intègre de point générique $t$, donc
à prouver la deuxième assertion de e). Celle-ci résulte aussitôt de d) et du
\margin{345}
\margin{345}
\margin{174}
